\chapter{Introdução}

Nenhuma norma é capaz de prever todas as situações em que será empregada. Quando se trata de uma área tão dinâmica e tão dependente da
tecnologia como a comunicação, a previsibilidade é ainda menor: novas mídias, linguagens e plataformas surgem a todo o momento. Por isso,
nesta introdução apresentaremos alguns princípios gerais que devem ser aplicados tanto a casos não previstos nesta \textbf{Política} quanto na
elaboração de novas versões deste documento.

\section{Uma política de comunicação pública}

No Brasil, a ideia de comunicação pública ganha relevância a partir da redemocratização, quando passa a substituir o conceito de ``comunicação
governamental'', introduzido em nossa administração pública durante os governos de Getúlio Vargas (1930--1945). O paradigma de comunicação
pública é, assim, uma tentativa de superar a comunicação unilateral, característica de governos autoritários, por meio da qual a população é
informada, mas não participa.

Para
\href{https://abcpublica.org.br/wp-content/uploads/2021/02/Comunica\%C3\%A7\%C3\%A3o-P\%C3\%BAblica-VF-Cap\%C3\%ADtulo.pdf}{Jorge
	Duarte}, ``fazer comunicação pública é assumir a perspectiva cidadã na comunicação envolvendo temas de interesse coletivo, alterando seu eixo,
tradicionalmente centrado no atendimento dos interesses da instituição e de seus gestores''. O objetivo da comunicação pública, conclui o autor,
é ``o atendimento do interesse público e da sociedade, simbolizado pelo cidadão''.

De maneira mais didática, a Associação Brasileira de Comunicação Pública (ABCPública) elenca \href{https://abcpublica.org.br/conheca-os-12-principios-da-comunicacao-publica/}{doze
	princípios da comunicação pública}:

\begin{enumerate}

	\item
	      garantir o acesso amplo à informação;
	\item
	      fomentar o diálogo;
	\item
	      estimular a participação;
	\item
	      promover os direitos e a democracia;
	\item
	      combater a desinformação;
	\item
	      ouvir a sociedade;
	\item
	      focar no cidadão;
	\item
	      ser inclusiva e plural;
	\item
	      tratar a comunicação como política de Estado;
	\item
	      garantir a impessoalidade;
	\item
	      pautar-se pela ética;
	\item
	      atuar com eficácia.
\end{enumerate}

\section{Princípios
  institucionais}

A comunicação do CRP~SP está pautada pelos princípios da ética, da responsabilidade pública, da transparência, da credibilidade e da colaboração. Materializa a voz e a \emph{persona} institucionais e, por meio dela, a autarquia mantém diálogo constante com a categoria e com a
sociedade. Empenha-se em comunicar a diversidade e a representatividade (racial, étnica, de gênero, de corpos, etária, de territórios). Seus produtos adotam linguagem inclusiva (simples, gendrada, antirracista, anticapacitista e com recursos de acessibilidade).

Os princípios institucionais a seguir estão presentes no \href{https://correio.crpsp.org.br/service/home/~/?auth=co&loc=pt_BR&id=4201&part=6}{Regimento
	Interno} do CRP~SP e são derivados das normativas referentes ao Sistema Conselhos de Psicologia e, principalmente, do \href{https://www.crpsp.org/uploads/pagina/289379/2j9LIMPLJ9jFr5YrK57HmAiBWjVMxdbe.pdf}{{Código de ética profissional da/o psicóloga/o}}. Além destes, publicações de caráter jornalístico estão sujeitas ao \href{https://fenaj.org.br/codigo-de-etica-dos-jornalistas-brasileiros/}{Código de ética dos jornalistas brasileiros}.

Também são documentos relacionados à \textbf{Política de Comunicação} e devem ser observados no desenvolvimento das atividades:

\begin{enumerate}

	\item
	      \href{https://www.planalto.gov.br/ccivil_03/constituicao/constituicao.htm}{Constituição federal};
	\item
	      \href{https://www.planalto.gov.br/ccivil_03/leis/l5766.htm}{Lei nº~5.766/71}, que criou o Sistema Conselhos de Psicologia;
	\item
	      Regimento interno do CRP~SP
	      (\href{https://transparencia.cfp.org.br/crp06/legislacao/regimento-interno-do-conselho-regional-de-psicologia-de-sao-paulo-6a-regiao-aprovado-pela-resolucao-no-5-de-22-de-marco-de-2023/}{Resolução CRP~SP nº~05/2023});
	\item
	      Código de ética profissional da/o psicóloga/o
	      (\href{https://atosoficiais.com.br/cfp/resolucao-do-exercicio-profissional-n-10-2005-aprova-o-codigo-de-etica-profissional-do-psicologo}{Resolução CFP n~º~010/2005});
	\item
	      \href{https://www.planalto.gov.br/ccivil_03/decreto/d70274.htm}{Decreto nº~70.274, de 9 de março de 1972} (normas do cerimonial público e a ordem geral de precedência);
	\item
	      Lei de acesso à informação (LAI; \href{https://www.planalto.gov.br/ccivil_03/_ato2011-2014/2011/lei/l12527.htm}{Lei nº~12.527/2011});
	\item
	      \href{https://dados.gov.br/dados/conteudo/politica-de-dados-abertos}{Política de dados abertos} do governo federal;
	\item 
		  \href{https://emag.governoeletronico.gov.br/}{{Modelo de acessibilidade em governo eletrônico}} (eMAG) do governo federal;
	\item
	      \href{https://transparencia.cfp.org.br/wp-content/uploads/sites/29/2024/08/Politica-de-TI-2-Edicao.pdf}{Política de segurança em tecnologia da informação e comunicação} do CRP~SP;
	\item
	      \href{https://transparencia.cfp.org.br/crp06/legislacao/portaria-crp-n-49-2024-institui-a-politica-de-patrocinio-e-apoio-institucional-do-conselho-regional-de-psicologia-da-6a-regiao-crp-06/}{Política de patrocínio e apoio institucional} do Conselho Regional de Psicologia da 6ª Região;
	\item
	      \href{http://cedoc.crpsp.org.br/handle/1/2938}{Manual de linguagem do CRP~SP}.
\end{enumerate}

\subsection{Respeito e promoção dos Direitos Humanos}

A centralidade da promoção dos Direitos Humanos e do respeito à dignidade da pessoa está expressa no Código de Ética Profissional da/o Psicóloga/o e no artigo 130 do
\href{https://correio.crpsp.org.br/service/home/~/?auth=co&loc=pt_BR&id=4201&part=6}{{Regimento interno}} do CRP~SP.

Além de prezar pela divulgação de pautas relacionadas ao tema, a comunicação da autarquia deve demonstrar seu compromisso com os Direitos
Humanos em todas as suas práticas. A linguagem adotada pelo Conselho, portanto, deve ser acima de tudo \textbf{inclusiva}: deve
respeitar as orientações da norma culta sem deixar de dar visibilidade a grupos social e historicamente marginalizados.

Assim, entendemos por linguagem inclusiva a linguagem \textbf{gendrada} (isto é, sem o uso do gênero masculino para generalizações),
\textbf{antirracista} (que não se utilize de palavras, expressões ou figuras de linguagem de conotação preconceituosa ou baseadas em
estereótipos de raça ou etnia) e \textbf{antietarista} (que não se utilize de palavras, expressões ou figuras de linguagem que tenham como
premissa uma visão negativa do envelhecimento).

A linguagem que adotamos também é \textbf{anticapacitista} --- não se utiliza de palavras, expressões ou figuras de linguagem que atribuam
valor a competências cognitivas, sensoriais ou motoras. Tão importante quanto usarmos uma linguagem que inclua as pessoas com deficiência,
contudo, é dispormos de \textbf{recursos de acessibilidade}: textos que sejam lidos adequadamente por leitores de tela, \emph{sites} de
navegação intuitiva, descrição de imagens e audiodescrição e uso de \textbf{linguagem simples}.

Nossos produtos são publicados em diversos formatos: no \emph{site}, como livro digital, na forma de postagens. Cada um desses formatos tem
diretrizes específicas de acessibilidade. Em qualquer caso, deverão ser seguidas as orientações do \href{https://emag.governoeletronico.gov.br/}{{Modelo de
				Acessibilidade em Governo Eletrônico}} (eMAG) do Governo Federal.

Como princípios fundamentais e complementares da comunicação do CRP~SP, \textbf{inclusão} e \textbf{acessibilidade} devem guiar toda e qualquer
produção da unidade de comunicação e todo e qualquer projeto de comunicação integrada.

\subsection{Transparência}

O artigo 131 do
\href{https://correio.crpsp.org.br/service/home/~/?auth=co&loc=pt_BR&id=4201&part=6}{{Regimento
			Interno}} do CRP~SP prevê a divulgação dos atos do Conselho com o
objetivo de aumentar o conhecimento e o reconhecimento da Psicologia
pela sociedade brasileira.

\subsection{Orientação}

O artigo 132 do
\href{https://correio.crpsp.org.br/service/home/~/?auth=co&loc=pt_BR&id=4201&part=6}{{Regimento
			Interno}} do CRP~SP determina que a autarquia mantenha ``publicações
destinadas à divulgação de matéria de interesse da psicóloga e do
público em geral, cabendo ao Plenário e/ou Diretoria dispor a respeito e
de acordo com a Política de Comunicação Institucional''.

\subsection{Efetivação de funções precípuas do Conselho Regional de
	Psicologia}

São funções precípuas dos conselhos profissionais, conforme o
\href{about:blank}{{Tribunal de Contas da União}} (TCU), cabendo à
Coordenação de Comunicação colaborar com seu desempenho:

\begin{enumerate}

	\item
	      orientar, supervisionar, fiscalizar e disciplinar o exercício da
	      profissão;
	\item
	      zelar pelo prestígio e bom nome da profissão e dos profissionais que a
	      exercem;
	\item
	      organizar e manter o registro profissional (registro, inscrição,
	      cancelamento);
	\item
	      examinar reclamações e representações;
	\item
	      representar às autoridades competentes;
	\item
	      atuar como órgão de consulta do governo;
	\item
	      atuar em colaboração com entidades de classe, escolas e faculdades;
	\item
	      contribuir para o desenvolvimento econômico;
	\item
	      disseminar a técnica econômica nos diversos setores da economia
	      nacional, promovendo estudos e campanhas em prol da racionalização
	      econômica do país;
	\item
	      promover estudos e campanhas em prol da racionalização administrativa
	      do país.
\end{enumerate}

\section{Princípios constitucionais aplicáveis à comunicação
  pública}

O
\href{https://www.planalto.gov.br/ccivil_03/decreto-lei/del0200.htm}{{Decreto-Lei
			nº~200}}, de 25 de fevereiro de 1967, define \textbf{autarquia} como

\begin{quote}
	{[}\ldots{]} o serviço autônomo, criado por lei, com personalidade
	jurídica, patrimônio e receita próprios, para executar atividades
	típicas da Administração Pública, que requeiram para seu melhor
	funcionamento, gestão administrativa e financeira descentralizada.
\end{quote}

Como autarquia, portanto, o CRP~SP está subordinado às normas e
princípios que regem a administração pública.

Por sua vez, o artigo 37 da
\href{https://www.planalto.gov.br/ccivil_03/constituicao/constituicao.htm}{{Constituição
			Federal}} prevê que ``A administração pública direta e indireta de
qualquer dos Poderes da União, dos Estados, do Distrito Federal e dos
Municípios obedecerá os princípios de legalidade, impessoalidade,
moralidade, publicidade e eficiência''. Esses princípios, recepcionados
no artigo 3º, inciso XXI do
\href{https://correio.crpsp.org.br/service/home/~/?auth=co&loc=pt_BR&id=4201&part=6}{{Regimento
			Interno do CRP~SP}}, são elucidados a seguir.

\subsection{Legalidade}

Diferentemente de instituições privadas, às quais é permitido fazer tudo
aquilo que não seja expressamente proibido em lei, às autarquias é
permitido fazer \textbf{apenas} o que a lei prevê.

A obediência ao princípio da legalidade deve, ainda, levar em conta a
hierarquia do ordenamento jurídico brasileiro, que tem a Constituição
Federal como lei maior, e abaixo dela as leis ordinárias --- das quais
destacamos a
\href{https://www.google.com/url?sa=t&source=web&rct=j&opi=89978449&url=https://www.planalto.gov.br/ccivil_03/leis/l5766.htm&ved=2ahUKEwi7voyM_MKJAxWYF7kGHcbCO54QFnoECBMQAw&usg=AOvVaw3BsZ1ugMxVxz5tlw9ZAIyv}{{Lei
			nº~5.766/1971}}, de 20 de dezembro de 1971, que cria o Conselho Federal
e os Conselhos Regionais de Psicologia, e a
\href{https://www.google.com/url?sa=t&source=web&rct=j&opi=89978449&url=https://www.planalto.gov.br/ccivil_03/leis/l9649cons.htm&ved=2ahUKEwjLibuZ_MKJAxVuILkGHaPnFtcQFnoECAoQAQ&usg=AOvVaw1dF6QYD5NfkI-TBa6TTJgo}{{Lei
			nº~9.649/1998}}, que dispõe, entre outras coisas, sobre os serviços de
fiscalização de profissões regulamentadas --- e as normativas próprias
da autarquia (portarias, resoluções etc.).

\subsection{Impessoalidade}

O princípio da impessoalidade determina que a atuação dos órgãos da
administração pública tem por finalidade atender ao interesse público.
No caso da comunicação do CRP~SP, isso implica tanto o atendimento das
demandas de profissionais de Psicologia quanto a defesa da categoria
perante a sociedade. A comunicação institucional da autarquia deve,
portanto, ser útil às psicólogas e aos psicólogos e responsável em sua
relação com a população em geral, pois representa a face oficial da
Psicologia no estado de São Paulo.

Como consequência do compromisso assumido com o interesse público, é
vedada a utilização da \emph{persona} institucional para a promoção de
indivíduos ou grupos. A comunicação institucional não pode fugir ao
princípio da legalidade, devendo restringir-se às situações previstas
nas normativas da instituição e na legislação pátria; segundo o § 1º~do
artigo 37 da
\href{https://www.google.com/url?sa=t&source=web&rct=j&opi=89978449&url=https://www.planalto.gov.br/ccivil_03/constituicao/constituicao.htm&ved=2ahUKEwjlis2B_MKJAxXcHrkGHTtTAvcQFnoECBgQAQ&usg=AOvVaw3i_8717crw9PBIV4q9Jndm}{{Constituição
			Federal}},

\begin{quote}
	A publicidade dos atos, programas, obras, serviços e campanhas dos
	órgãos públicos deverá ter caráter educativo, informativo ou de
	orientação social, dela não podendo constar nomes, símbolos ou imagens
	que caracterizem promoção pessoal de autoridades ou servidores públicos.
\end{quote}

\subsection{Moralidade}

O conceito administrativo de moralidade está ligado às ideias de
probidade, decoro e boa-fé. Nesse sentido, ``moralidade'' não é um
conceito subjetivo, mas um referencial de atuação adequada aos
pressupostos éticos adotados por um dado grupo social. Na administração
pública, esses pressupostos são definidos pela legislação que define os
limites de atuação de uma agente pública.

\subsection{Publicidade}

Como princípio da administração pública, a publicidade se refere tanto à
exigência de publicação de atos administrativos para que produzam
efeitos externos --- a publicação de uma portaria, por exemplo ---
quanto à exigência de transparência da atuação administrativa. É
necessário, por exemplo, que os atos administrativos apresentem sua
motivação, isto é, que explicitem os pressupostos fáticos e jurídicos
que o determinam, para que sociedade e órgãos de controle possam
fiscalizar a legitimidade desses atos.

\subsection{Eficiência}

Por fim, o princípio de eficiência fala tanto da qualidade da atuação da
agente pública quanto da racionalidade da administração. A eficiência
das agentes é determinada pela qualidade e produtividade de seu
trabalho; a racionalidade da administração está ligada à prestação de
serviço público adequado, definido pela
\href{https://www.planalto.gov.br/ccivil_03/leis/l8987cons.htm}{Lei nº~8.987/95} como aquele que ``satisfaz as condições de regularidade,
continuidade, eficiência, segurança, atualidade, generalidade, cortesia
na sua prestação e modicidade das tarifas'' (art. 6º, § 1º).

\section{Implementação e gestão desta
  Política}

Esta \textbf{Política} deverá ser observada por conselheiras e
conselheiros, bem como por trabalhadoras e trabalhadores do CRP~SP.

Questões não previstas nesta \textbf{Política} deverão ser encaminhadas
para o \emph{e-mail}
\href{mailto:ascom@crpsp.org.br}{{ascom@crpsp.org.br}}, para
avaliação e proposição às instâncias competentes.